\begin{helper}[An empty eligible family locates common missing colors]
Fix a finite loopless multigraph with individual edge identities,
a palette of $k$ colors, a distinguished edge $e_0$, a certified input
coloring, an exact ordered tree $T$, a generating coloring $\varphi$,
an ETT $\mathcal E$, a \noderef{HierarchyBase}, a
\noderef{HierarchyCounterexample} $C$, and its
\noderef{HierarchyMinimumCounterexampleContext}.
Write $A=\noderef{HierarchyStarredStage}(C.\mathit{hierarchy},q)$ at
the last stored index, $\ell=\noderef{HierarchyFullTailLength}(C)$,
and $j=\noderef{HierarchyPathNumber}(C.\mathit{hierarchy})$.
Assume $2\le j<\ell$ and that the original-relative finite family
\noderef{HierarchyEligibleCommonMissingIndices} is empty.
Let $P$ be any original-relative
\noderef{HierarchyReplacementPresentation} and let $c$ be missing
under its coloring both on the original exact penultimate prefix
$B=\noderef{HierarchyFullTailPrefix}(C,\ell-1)$ and at the last
vertex of $T$. Then $c$ is missing on the original exact $A$ under
$P.\mathit{coloring}$.
\end{helper}
\begin{proof}
The fields of \noderef{TashkinovHierarchy} give
$A\preccurlyeq T$: at divider zero use the base core's inclusion
in the first stage, and otherwise use the final nesting inclusion.
The definition of \noderef{HierarchyFullTailLength} gives
$\ell=|T|-|A|$. Hence $\ell\ge3$ and the definition of
\noderef{HierarchyFullTailPrefix} gives $|B|=|T|-1$.
Both $A$ and $B$ are exact ordered prefixes of $T$, with
$|A|\le|B|$, so $A\preccurlyeq B$: at every vertex or edge index
the corresponding entries are equal to the same entry of $T$.
These are all clauses of \noderef{TreeSequenceInitialSegment}.

Expand \noderef{TreeMissingColors} and
\noderef{TreeSequenceVertexSet} to choose an index
$x\in\{0,\ldots,|B|\}$ with $c$ missing at $B.\mathit{vertex}(x)$.
If $x\le|A|$, the exact-prefix vertex equality puts that same
witness in $A$, proving the conclusion.

Otherwise put $i=x-|A|$. Then $1\le i\le\ell-1$ and the exact
prefix $T_i=\noderef{HierarchyFullTailPrefix}(C,i)$ has length
$|A|+i=x$, since $x\le|B|<|T|$. Its last vertex is exactly the
displayed vertex at index $x$ of $T$, hence exactly
$B.\mathit{vertex}(x)$. Thus $c$ is missing under $P.\mathit{coloring}$
at both the last vertex of $T_i$ and the last vertex of $T$.
The bounds on $i$, the supplied original-relative $P$, and these
two missing witnesses are precisely the membership conditions in
the filter defining \noderef{HierarchyEligibleCommonMissingIndices}.
That family therefore contains $i$, contradicting its emptiness.
This rules out the second alternative and proves the assertion.
\end{proof}
